\documentclass[10pt,a4paper]{amsart}
\usepackage[margin=19mm]{geometry}
\usepackage{amsmath,amssymb,mathtools}
\usepackage[scr=rsfs,cal=euler]{mathalfa}
\usepackage{xfrac}
\usepackage[unicode,hidelinks,bookmarks=false]{hyperref}
\newcommand{\ie}{i.e.\ }
\newcommand{\cf}{cf.\ }
\newcommand{\resp}{respectively\ }
\newcommand{\Subsection}{\subsection}
\providecommand{\nobreakdash}{\nobreak\hbox{-}}
\setlength{\emergencystretch}{2em}
\multlinegap0pt
\usepackage{bm}
\usepackage{bbm}
\usepackage{dsfont}
\usepackage{xspace}
\usepackage{cancel}
\usepackage{mathtools}
\usepackage{amsthm}
\makeatletter
\@ifundefined{theorem}{\newtheorem{theorem}{Theorem}}{}
\@ifundefined{proposition}{\newtheorem{proposition}[theorem]{Proposition}}{}
\makeatother
\title[Directional Subdifferentials of the Value Function in Asplun]{Directional Subdifferentials of the Value Function in Asplund Spaces}
\author{Weihao Mao, Jane J. Ye}
\date{Source: arXiv:2608.20241v1 (CC BY 4.0)}
\begin{document}
\maketitle

\noindent\emph{Source and licence.} Directional Subdifferentials of the Value Function in Asplund Spaces. Version arXiv:2608.20241v1 (CC BY 4.0), distributed under the Creative Commons Attribution 4.0 International licence, as recorded in the arXiv licence metadata for this exact version and shown on the abstract page https://arxiv.org/abs/2608.20241v1. Full rights notice: licenses/arxiv-2608.20241-CC-BY-4.0.txt.\par\smallskip

\noindent\textit{Editorial scope.} Counted unit: the Theorem whose source label is \texttt{thm: chain rule} in the article (source0.tex lines 469--489, source label \texttt{thm: chain rule}), delivered together with the article's own complete original proof of it (source0.tex lines 491--553, one \texttt{proof} environment of 63 lines). No mathematics is written, repaired or re-derived: statement and proof are the article's own text line for line. Required context retained uncounted: prop of directional normal cones of product of separated sets (lines 166--174); thm: directional normal cones of preimage sets (lines 325--334); thm:scalarization (lines 435--440). Every definition, display and label of those lines is retained unchanged. Omitted from this package and not redistributed: 1--165, 175--324, 335--434, 441--468, 490--490, 554--1212. Nothing inside a retained excerpt is omitted or altered: every retained body is delivered exactly as the article's lines are written, so no omission receipt of any kind accompanies this package. Citations: the retained excerpts carry no citation command at all, so no bibliography entry of the article is transcribed and no reference is redistributed. Reference adaptation: none was needed and none was made; every reference inside the delivered unit resolves inside it, so no unresolved reference or citation is produced. Printed slips kept literal: no printed word, symbol, hypothesis or number of the counted statement or of its proof was altered. Layout and class: the article's own class is \texttt{amsart} with packages including mathptmx, amsmath, amssymb, amsfonts, amsthm, enumerate, mathrsfs, xcolor, hyperref; the wrapper is the lane's pinned \texttt{amsart} editorial wrapper at 10pt on A4, which restates the theorem-like environments the retained text needs and adds the article's own macro definitions that the retained text needs (none), with extra wrapper packages (bm, bbm, dsfont, xspace, cancel, mathtools). The delivered numbering is the unit's own. Assets: no retained span contains a figure environment, a graphics-inclusion command, a PDF-page inclusion, a table or a TikZ picture; no image file is copied into this package, the package declares no asset dependency and no image file is redistributed with it. Licence: version arXiv:2608.20241v1 of this article is distributed under the Creative Commons Attribution 4.0 International licence, as recorded in the arXiv licence metadata for this exact version and shown on the abstract page https://arxiv.org/abs/2608.20241v1. A full rights notice is delivered at licenses/arxiv-2608.20241-CC-BY-4.0.txt. Characters: every retained excerpt is pure ASCII, so the delivered engine, XeLaTeX with its default Latin Modern text font, reports no missing character.
\begin{proposition}\label{prop of directional normal cones of product of separated sets}
Let $\Omega_i\subset X_i$ $(i=1,2)$ be nonempty subsets, and let $\bar x=(\bar x_1,\bar x_2)\in \Omega:=\Omega_1\times \Omega_2$. 
Then, for any $\bar u=(\bar u_1,\bar u_2)\in X_1\times X_2$, one has
\begin{equation}
    N(\bar x;\Omega;\bar u)
    \subset 
    N(\bar x_1;\Omega_1;\bar u_1)\times N(\bar x_2;\Omega_2;\bar u_2).
\end{equation}
\end{proposition}

\begin{theorem}\label{thm: directional normal cones of preimage sets} 
Let \( X \) and \( Z \) be Asplund spaces, and \( Q \subset Z \) be closed.  
Consider a continuous mapping \( \phi : X \to Z \), and define
$C = \phi^{-1}(Q):= \{x\mid \phi(x)\in Q\}, ~F(x) := Q - \phi(x)$. 
Assume that \( F \) is metrically subregular at \( (\bar{x}, 0) \) in some direction \( h \in X \).  
Suppose further that \( \phi \) is Hadamard differentiable at \(\bar{x}\) in direction \( h \) with directional derivative $v := \phi'_H(\bar{x}; h)\in T(\phi(\bar x);Q)$. Then the following inclusion holds:
\begin{equation}\label{eq:dir_normal_preimage}
    N(\bar{x}; C; h) \;\subset\; D^*_N \phi(\bar{x}; h) \big( N(\phi(\bar{x}); Q; v) \big).
\end{equation}
\end{theorem}

\begin{theorem} \label{thm:scalarization}
Let $X,Y$ be Asplund spaces. Suppose that the mapping $\phi: X \to Y$ is locally Lipschitzian and weak* strictly Lipschitzian at $\bar{x} \in X$. If $\phi$ is Hadamard differentiable at $\bar{x}$ in a given direction $\bar{u} \in X$, then for any $y^* \in Y^*$, we have: 
\begin{equation} \label{eq:scalarization_formula}
    D_N^* \phi(\bar{x}; \bar{u})(y^*) = \partial \langle y^*, \phi \rangle (\bar{x}; \bar{u}).
\end{equation}
\end{theorem}

\begin{theorem}[Chain Rule]\label{thm: chain rule}
Let \( X, Z \) be Asplund spaces, and let \( \phi : X \to Z \) be continuous.  
Let \( g : Z \to \overline{\mathbb{R}} \) be finite at \( \phi(\bar{x}) \) and lower semicontinuous, and define \( \psi := g \circ \phi \).  
Given a direction \( (h, 0) \in X \times \mathbb{R} \), suppose the mapping
\[
    F : X \times \mathbb{R} \rightrightarrows Z \times \mathbb{R}, 
    \quad F(x, \alpha) := \operatorname{epi} g - (\phi(x), \alpha),
\]
is metrically subregular at \( ((\bar{x}, \psi(\bar{x})), (0,0)) \) in direction \( (h, 0) \), and $\phi$ is Hadamard differentiable in $h$ at $\bar x$ with \( v = \phi'_H(\bar{x}; h) \).  
Then
\begin{equation}\label{eq:chain_rule_unscalarized}
    \partial \psi(\bar{x}; h)
    \;\subset\;
    D^*_N \phi(\bar{x}; h) \big( \partial g(\phi(\bar{x}); v)\big).
\end{equation}
Furthermore, if $\phi$ is weak* strictly Lipschitzian at $\bar{x}$, then 
\begin{equation}\label{eq:chain_rule}
    \partial \psi(\bar{x}; h)
    \;\subset\;\bigcup\limits_{z^*\in \partial g(\phi(\bar{x}); v)} \partial \langle z^*, \phi\rangle(\bar{x}; h).
\end{equation}
\end{theorem}

\begin{proof}
Take any $x^*\in \partial \psi(\bar x; h)$. By the definition of the directional limiting subdifferential, and for the epi-direction $(h, 0)$, we have  
\( (x^*, -1) \in N((\bar{x}, \psi(\bar{x})); \operatorname{epi} \psi; (h, 0)) \).  
For any \( (x, \alpha) \in \operatorname{epi} \psi \), we have a bijective correspondence with \( (\phi(x), \alpha) \), so that
\[
    \operatorname{epi} \psi = \Psi^{-1}(\operatorname{epi} g),
    \quad \text{where } \Psi(x, \alpha) := (\phi(x), \alpha).
\]
Note that the mapping $\Psi$ is Hadamard differentiable at $(\bar{x}, \psi(\bar{x}))$ in the direction $(h, 0)$, yielding $\Psi'_H(\bar{x}, \psi(\bar{x}); (h, 0)) = (\phi'_H(\bar{x}; h), 0) = (v, 0)$.
Applying Theorem~\ref{thm: directional normal cones of preimage sets} with \( \Psi \) in place of \( \phi \), \( \operatorname{epi} g \) in place of \( Q \) and $C:=\{(x,\alpha) \mid (\phi(x),\alpha)\in \text{epi}g)\}$ therein,   
we deduce that
\begin{equation}\label{eq:preimage_chain}
    (x^*, -1) 
    \in D^*_N \Psi((\bar{x},\psi(\bar x)); (h, 0)) \,
    \big( N((\phi(\bar{x}), \psi(\bar{x})); \operatorname{epi} g; (v, 0)) \big).
\end{equation}

Hence, there exist multipliers \( (y^*, \beta) \) such that
\begin{equation}\label{eq:normals_gphPsi}
    (x^*, -1, -y^*, -\beta)
    \in N\bigl((\bar{x}, \psi(\bar{x}), \phi(\bar{x}), \psi(\bar{x}));
               \operatorname{gph} \Psi; (h, 0, v, 0)\bigr).
\end{equation}
Define $\Omega_\Psi
    := \{ (a, b, c, d) \in X \times Z \times \mathbb{R}^2 
        \mid (a, c, b, d) \in \operatorname{gph} \Psi \}$. 
Then, by the definition of the directional normal cone, condition~\eqref{eq:normals_gphPsi} is equivalent to
\begin{equation}\label{eq:normals_omega}
    (x^*, -y^*, -1, -\beta)
    \in N\bigl((\bar{x}, \phi(\bar{x}), \psi(\bar{x}), \psi(\bar{x}));
               \Omega_\Psi; (h, v, 0, 0)\bigr).
\end{equation}
By Proposition~\ref{prop of directional normal cones of product of separated sets},  
and noting that $\Omega_\Psi = \operatorname{gph} \phi \times \operatorname{gph} \operatorname{Id}|_{\psi}$, 
where \( \operatorname{Id}|_{\psi} \) denotes the identity mapping restricted to the range of \( \psi \),  
we can decompose the directional normal cone:
\begin{equation}\label{eq:decomposition of normal cone in chain rule}
    \begin{aligned}
    &N\bigl((\bar{x}, \phi(\bar{x}), \psi(\bar{x}), \psi(\bar{x}));
        \operatorname{gph} \phi \times 
        \operatorname{gph} \operatorname{Id}|_{\psi}; 
        (h, v, 0, 0)\bigr) \\
    &\subset
    N((\bar{x}, \phi(\bar{x})); \operatorname{gph} \phi; (h, v))
    \times
    N((\psi(\bar{x}), \psi(\bar{x})); 
      \operatorname{gph} \operatorname{Id}|_{\psi}; (0, 0)).
    \end{aligned}
\end{equation}
Thus we have
\begin{equation}
    \begin{split}
        &(x^*, -y^*) \in N((\bar{x}, \phi(\bar{x})); \operatorname{gph} \phi; (h, v)),\\
&(-1, -\beta) \in 
N((\psi(\bar{x}), \psi(\bar{x})); 
  \operatorname{gph} \operatorname{Id}|_{\psi}; (0, 0)). 
    \end{split}
\end{equation}
Because the direction is identically zero, the latter condition is evaluated on the standard normal cone to the identity graph, which strictly enforces $-1 - \beta = 0$, leading to \( \beta = -1 \).  
Consequently, $(y^*, -1) \in N((\phi(\bar{x}), \psi(\bar{x})); \operatorname{epi} g; (v, 0))$, meaning \( y^* \in \partial g(\phi(\bar{x}); v) \). It follows that $x^* \in D^*_N \phi(\bar{x}; h)(y^*)$. 
This proves the unscalarized inclusion \eqref{eq:chain_rule_unscalarized}. 
The scalarized inclusion \eqref{eq:chain_rule} then follows immediately by Theorem \ref{thm:scalarization}. 
\end{proof}

\end{document}
